Una dieta debe aportar al menos 80 unidades de proteína usando dos alimentos
$x_1$ y $x_2$, sin superar 60 unidades de grasa. Cada unidad de $x_1$ aporta 2 de
proteína y 4 de grasa; cada unidad de $x_2$, 4 de proteína y 3 de grasa. Los
costes unitarios son 3 y 2. El modelo que minimiza el coste es:

\begin{equation}
\begin{split}
\label{eq:demo_dieta}
\mbox{min. } z = 3x_1 + 2x_2\\
s.a.:\\
2x_1 + 4x_2 \geq 80\\
4x_1 + 3x_2 \leq 60\\
x_1,\,\,x_2\geq 0\\
\end{split}
\end{equation}

\begin{enumerate}
    \item Resolver el problema mediante el método de Lemke.
\end{enumerate}
